\section{Algorithmic Composition}\label{sec:background-algorithmic-composition}

Algorithmic composition is the technique of using algorithms—formal sets of rules or instructions—to create music.
While this concept has roots in classical music (such as Mozart's musical dice games), the advent of
computers has exponentially expanded its scope and capabilities.

\subsection{Procedural Generation vs. Direct Input}

Traditional digital music production relies on Direct Input.
A composer uses a MIDI keyboard or a graphical Piano Roll in a Digital Audio Workstation (DAW) to manually
input every note, chord, and rest.
This process grants total control over the micro-details of a performance but can be tedious when
constructing complex, repetitive, or mathematically structured sections.

Algorithmic composition shifts the paradigm to Procedural Generation.
Instead of writing the notes, the composer writes the \textit{rules} that generate the notes.
For example, rather than manually copying and pasting a bassline across 32 bars while applying slight
variations to the velocity, a composer might write a simple loop that algorithmically modifies the velocity
parameter on each iteration.

\subsection{Structured Abstraction}

The true power of algorithmic composition lies in its ability to abstract musical ideas.
A chord progression can be represented as an array of data points, and a rhythm can be defined as a
mathematical function over time.
This approach allows composers to experiment rapidly: changing a foundational variable (such as a root note
or a time signature) can instantly cascade throughout the entire algorithmic structure, generating a
completely new variation of the piece without the need for manual editing.
The DSL proposed in this thesis aims to provide a robust, type-safe environment specifically optimized for
this methodology of structured, algorithmic creation.
